% GENERATED FILE. Edit canonical lessons, then run npm run generate:books.
% canonical-source-hash: fnv1a64:db929978755a799e
% canonical-lessons: FR-C436-raccourci, FR-C436-detour, FR-C436-itineraire, FR-C436-contourner, FR-C436-rebrousser

\chapter{Raccourcis et détours}
\label{ch:fr-raccourcis-et-detours}
\hlchaptermodality{436}

\begin{chapteropening}
\textbf{By the end of this chapter:} \emph{I can talk about routes: the shortcut, the detour, the route plan, going round something and turning back.}

Complete the last lesson of chapter 436: I can talk about routes: the shortcut, the detour, the route plan, going round something and turning back.
\end{chapteropening}

\section[le raccourci]{le raccourci — the shortcut}

\label{lesson:FR-C436-raccourci}



\begin{quote}
Before the new one: say the French for the marker post, then the French for the corner.
\end{quote}

\subsection*{You'll want to know: le raccourci}
\textbf{le raccourci} — \textquotedblleft{}the shortcut\textquotedblright{}.

\begin{quote}
Il connaît un raccourci par le parc.
\end{quote}

\begin{quote}
\emph{He knows a shortcut through the park.}
\end{quote}

Built on \emph{raccourcir} 'to shorten'.

\begin{cousinweb}
\emph{raccourci} is the past participle of \emph{raccourcir} 'to shorten', used as a noun (EtymDB); the verb is built on \emph{court} 'short', from Latin \emph{curtus} (Lewis \& Short).
\end{cousinweb}

\begin{grammarlens}[title={What you've built}]
One more word for telling, explaining and describing at B1.
\end{grammarlens}

\subsection*{Your turn}
\begin{itemize}
\raggedright
  \item \emph{Say it:} \emph{le raccourci}
  \item \emph{Say it:} \emph{le raccourci}, once more
  \item \emph{Say it:} a sentence of your own with \emph{le raccourci}
  \item \emph{Recall:} say the French for the marker post, then the French for the corner, then say \emph{le raccourci} again
\end{itemize}

\begin{tcolorbox}[breakable,colback=teal!4,colframe=teal!35!black,title={Before you move on}]
What does le raccourci mean? (\textquotedblleft{}The shortcut\textquotedblright{}.) Say it once more.
\end{tcolorbox}
\section[le détour]{le détour — the detour, the roundabout way}

\label{lesson:FR-C436-detour}



\begin{quote}
Before the new one: say the French for the corner, then the French for the shortcut.
\end{quote}

\subsection*{You'll want to know: le détour}
\textbf{le détour} — \textquotedblleft{}the detour, the roundabout way\textquotedblright{}.

\begin{quote}
Nous avons fait un détour pour voir le château.
\end{quote}

\begin{quote}
\emph{We made a detour to visit the castle.}
\end{quote}

\emph{Sans détour} — straight out, frankly.

\begin{cousinweb}
\emph{détour} comes from Old French \emph{destour}, built on \emph{détourner} 'to turn away' (EtymDB), from \emph{tourner}, which goes back to Latin \emph{tornāre} 'to turn on a lathe'.
\end{cousinweb}

\begin{grammarlens}[title={What you've built}]
One more word for telling, explaining and describing at B1.
\end{grammarlens}

\subsection*{Your turn}
\begin{itemize}
\raggedright
  \item \emph{Say it:} \emph{le détour}
  \item \emph{Say it:} \emph{le détour}, once more
  \item \emph{Say it:} a sentence of your own with \emph{le détour}
  \item \emph{Recall:} say the French for the corner, then the French for the shortcut, then say \emph{le détour} again
\end{itemize}

\begin{tcolorbox}[breakable,colback=teal!4,colframe=teal!35!black,title={Before you move on}]
What does le détour mean? (\textquotedblleft{}The detour, the roundabout way\textquotedblright{}.) Say it once more.
\end{tcolorbox}
\section[l'itinéraire]{l'itinéraire — the route, the itinerary}

\label{lesson:FR-C436-itineraire}



\begin{quote}
Before the new one: say the French for the shortcut, then the French for the detour.
\end{quote}

\subsection*{You'll want to know: l'itinéraire}
\textbf{l'itinéraire} — \textquotedblleft{}the route, the itinerary\textquotedblright{}.

\begin{quote}
Avant de partir, elle prépare son itinéraire.
\end{quote}

\begin{quote}
\emph{Before leaving, she plans her route.}
\end{quote}

A masculine noun. \emph{Un itinéraire bis} — an alternative route.

\begin{cousinweb}
\emph{itinéraire} is a learned borrowing from Late Latin \emph{itinerārium} 'travel guide' (EtymDB), from \emph{iter, itineris} 'journey'.
\end{cousinweb}

\begin{grammarlens}[title={What you've built}]
One more word for telling, explaining and describing at B1.
\end{grammarlens}

\subsection*{Your turn}
\begin{itemize}
\raggedright
  \item \emph{Say it:} \emph{l'itinéraire}
  \item \emph{Say it:} \emph{l'itinéraire}, once more
  \item \emph{Say it:} a sentence of your own with \emph{l'itinéraire}
  \item \emph{Recall:} say the French for the shortcut, then the French for the detour, then say \emph{l'itinéraire} again
\end{itemize}

\begin{tcolorbox}[breakable,colback=teal!4,colframe=teal!35!black,title={Before you move on}]
What does l'itinéraire mean? (\textquotedblleft{}The route, the itinerary\textquotedblright{}.) Say it once more.
\end{tcolorbox}
\section[contourner]{contourner — to go round, to bypass}

\label{lesson:FR-C436-contourner}



\begin{quote}
Before the new one: say the French for the detour, then the French for the route.
\end{quote}

\subsection*{You'll want to know: contourner}
\textbf{contourner} — \textquotedblleft{}to go round, to bypass\textquotedblright{}.

\begin{quote}
La route contourne la ville par le nord.
\end{quote}

\begin{quote}
\emph{The road goes round the town to the north.}
\end{quote}

Figurative: \emph{contourner un problème} — to get round a problem.

\begin{cousinweb}
\emph{contourner} goes back to Vulgar Latin \emph{contornāre} (EtymDB), from Latin \emph{tornāre} 'to turn'.
\end{cousinweb}

\begin{grammarlens}[title={What you've built}]
One more word for telling, explaining and describing at B1.
\end{grammarlens}

\subsection*{Your turn}
\begin{itemize}
\raggedright
  \item \emph{Say it:} \emph{contourner}
  \item \emph{Say it:} \emph{contourner}, once more
  \item \emph{Say it:} a sentence of your own with \emph{contourner}
  \item \emph{Recall:} say the French for the detour, then the French for the route, then say \emph{contourner} again
\end{itemize}

\begin{tcolorbox}[breakable,colback=teal!4,colframe=teal!35!black,title={Before you move on}]
What does contourner mean? (\textquotedblleft{}To go round, to bypass\textquotedblright{}.) Say it once more.
\end{tcolorbox}
\section[rebrousser]{rebrousser — to turn back (rebrousser chemin)}

\label{lesson:FR-C436-rebrousser}



\begin{quote}
Before the new one: say the French for the route, then the French for to go round.
\end{quote}

\subsection*{You'll want to know: rebrousser}
\textbf{rebrousser} — \textquotedblleft{}to turn back (rebrousser chemin)\textquotedblright{}.

\begin{quote}
La route était fermée, alors nous avons rebroussé chemin.
\end{quote}

\begin{quote}
\emph{The road was closed, so we turned back.}
\end{quote}

Almost always in \emph{rebrousser chemin} — to retrace one's steps.

\begin{cousinweb}
\emph{rebrousser} is usually traced to \emph{rebours} 'against the grain', from Late Latin \emph{reburrus} 'with hair standing up'; the details are not settled, and EtymDB does not list it.
\end{cousinweb}

\begin{grammarlens}[title={What you've built}]
One more word for telling, explaining and describing at B1.
\end{grammarlens}

\subsection*{Your turn}
\begin{itemize}
\raggedright
  \item \emph{Say it:} \emph{rebrousser}
  \item \emph{Say it:} \emph{rebrousser}, once more
  \item \emph{Say it:} a sentence of your own with \emph{rebrousser}
  \item \emph{Recall:} say the French for the route, then the French for to go round, then say \emph{rebrousser} again
\end{itemize}

\begin{tcolorbox}[breakable,colback=teal!4,colframe=teal!35!black,title={Before you move on}]
What does rebrousser mean? (\textquotedblleft{}To turn back (rebrousser chemin)\textquotedblright{}.) Say it once more.
\end{tcolorbox}
